\subsection{\problemblock{*Iip} Data Block}\label{sec:block-iip}

The initial impact point (IIP) is obtained by integrating a ballistic trajectory from the
current vehicle state to a final altitude with no propulsive forces, as described in
\cref{sec:iip-calculations}. A simplified constant-drag aerodynamic model is
specified by a ballistic coefficient. The result estimates the impact point following a
propulsive-system failure and is used in range-safety analysis.

The calculation is performed when any associated output variable is referenced:
\statevar{iip_latgd}, \statevar{iip_long}, \statevar{iip_time},
\statevar{iip_rng}, or \statevar{iip_azm}. The required ballistic coefficient and
final impact altitude are given by

\begin{lstlisting}[style=taosmanual]
*iip iip_beta=550 iip_alt=1000
\end{lstlisting}

which sets the coefficient to 550~lb/ft$^2$ and the impact altitude to 1000~ft when
default units are used. Both values default to zero. A zero-drag calculation is
requested by setting the ballistic coefficient to zero rather than to a very large
number. The coefficient may be reset or incremented at any time during the
trajectory.

The outputs \statevar{iip_rng} and \statevar{iip_azm} give range and azimuth from
the origin of the tangent-plane coordinate system. When those quantities are used,
the \problemblock{*tangent} block should define the desired reference origin.
